% !TeX program = xelatex
% =====================================================================
%  2026 年全国大学生数学建模竞赛 B 题 —— 论文主文件（电子版）
%
%  依据《全国大学生数学建模竞赛论文格式规范（2026年修订稿）》：
%   - 电子版论文第一页必须为「摘要专用页」；
%   - 承诺书、编号专用页不得出现在电子版中（见规范第十条）。
%  故采用 withoutpreface 选项：去除承诺书与编号专用页，首页即摘要页。
%
%  编译方式（必须使用 xelatex）：
%      latexmk -xelatex main.tex
%   或  xelatex main.tex （连续运行两次以更新交叉引用）
% =====================================================================
\documentclass[withoutpreface]{cumcmthesis}

% =====================================================================
% 排版紧凑化（适中档）：覆盖模板默认的宽松参数，使整体更紧凑。
%   - 正文行距 1.38 -> 1.25
%   - 三线表行距 1.38 -> 1.0
%   - 列表顶部间距收紧（topsep -> 0）
%   - 章节标题前后距收紧
%   - 行间公式上下距收紧
%  如需恢复模板原样，删除本块即可。
% =====================================================================
\makeatletter
% 关键词：2026 规范为「关键词」（模板默认「关键字」）
\renewcommand*{\mcm@cap@keywordsname}{关键词}
% 正文行距
\renewcommand*{\baselinestretch}{1.25}
% 三线表行距（复用模板保存的原始 tabular）
\renewenvironment{tabular}%
  {\bgroup\renewcommand{\arraystretch}{1.0}\mcm@oldtabular}%
  {\mcm@endoldtabular\egroup}
% 章节标题前后距（前距 / 后距，单位 ex；数值越小越紧凑）
\renewcommand\section{\@startsection{section}{1}{\z@}%
  {-2ex \@plus -1ex \@minus -.2ex}%
  {0.8ex \@plus.2ex}%
  {\centering\normalfont\Large\bfseries}}
\renewcommand\subsection{\@startsection{subsection}{2}{\z@}%
  {-1.5ex\@plus -1ex \@minus -.2ex}%
  {0.5ex \@plus .2ex}%
  {\normalfont\large\bfseries}}
\renewcommand\subsubsection{\@startsection{subsubsection}{3}{\z@}%
  {-1.5ex\@plus -1ex \@minus -.2ex}%
  {0.5ex \@plus .2ex}%
  {\normalfont\normalsize\bfseries}}
% 参考文献：收紧条目间距（thebibliography 用原生 \list，不受上面 \setlist 影响）
\renewenvironment{thebibliography}[1]
     {\section*{\refname}%
      \@mkboth{\MakeUppercase\refname}{\MakeUppercase\refname}%
      \list{\@biblabel{\@arabic\c@enumiv}}%
           {\settowidth\labelwidth{\@biblabel{#1}}%
            \leftmargin\labelwidth
            \advance\leftmargin\labelsep
            \itemsep=0pt\parsep=0pt\topsep=0pt
            \@openbib@code
            \usecounter{enumiv}%
            \let\p@enumiv\@empty
            \renewcommand\theenumiv{\@arabic\c@enumiv}}%
      \sloppy
      \clubpenalty4000
      \@clubpenalty \clubpenalty
      \widowpenalty4000%
      \sfcode`\.\@m}
     {\def\@noitemerr
       {\@latex@warning{Empty `thebibliography' environment}}%
      \endlist}
\makeatother

% 列表间距（enumitem）：顶部间距归零，其余保持紧凑
\setlist{topsep=0em, partopsep=0pt, itemsep=0ex plus 0.1ex, parsep=0pt,
  leftmargin=1.5em, rightmargin=0em, labelsep=0.5em, labelwidth=2em}

% 行间公式上下距（须在 \begin{document} 之后设置，因 \normalsize 会重置这些长度）
\AtBeginDocument{%
  \setlength{\abovedisplayskip}{8pt plus 2pt minus 4pt}%
  \setlength{\belowdisplayskip}{8pt plus 2pt minus 4pt}%
  \setlength{\abovedisplayshortskip}{0pt plus 2pt}%
  \setlength{\belowdisplayshortskip}{4pt plus 2pt minus 2pt}%
}

% 图表标题（caption）上下间距收紧：caption 包默认 above=10pt 偏大
\setlength{\abovecaptionskip}{4pt}% 图题↔图片 / 表题↔正文
\setlength{\belowcaptionskip}{2pt}% 图题↔正文 / 表题↔表格

% 浮动体（图/表整体）与正文、浮动体之间的上下间距收紧
% 默认 textfloatsep=20pt、floatsep=12pt、intextsep=12pt，偏大
\setlength{\textfloatsep}{8pt plus 2pt minus 2pt}% 正文↔页顶/页底浮动体
\setlength{\floatsep}{6pt plus 2pt minus 2pt}%   相邻浮动体之间
\setlength{\intextsep}{6pt plus 2pt minus 2pt}%  正文中[h]浮动体↔上下正文

% ---------------------- 论文信息（按需填写） ----------------------
\title{无线电干扰源的快速自动定位与清除模型} % 【标题可待定，摘要页居中显示】
\tihao{B}                 % 题号：B
\baominghao{}             % 报名队号（电子版不打印，留空即可）
\schoolname{}             % 参赛学校（电子版严禁出现，留空）
\membera{}                % 队员 1（电子版不打印）
\memberb{}                % 队员 2（电子版不打印）
\memberc{}                % 队员 3（电子版不打印）
\supervisor{}             % 指导教师（电子版不打印）
\yearinput{2026}
\monthinput{09}
\dayinput{13}

\begin{document}

 \maketitle

 % ---------------------- 摘要专用页（不超过一页） ----------------------
 \begin{abstract}
 % =====================================================================
%  摘要 + 关键词（占位，最终成文时替换）
%  要求：概括全文针对 B 题四个子问题建立的模型、方法与主要结果；
%        无需翻译成英文；原则上不超过一页。
% =====================================================================
【待填充：全文摘要，建议 300--500 字。写作提纲（可按此展开）：
  (1) 背景与目标——无线电干扰源的快速自动定位与清除；
  (2) 针对问题一，建立示向度交会定位模型，……；
  (3) 针对问题二，建立第二检测点最优选择模型，……；
  (4) 针对问题三，建立多源快速定位与清除策略，……；
  (5) 针对问题四，建立定向干扰源定位模型，……；
  (6) 主要结论与数值/实测指标（清除比例、平均定位清除时间等）。】

本文针对无线电干扰源的快速自动定位与清除问题，\ldots{}（此处填写全文摘要）。

\keywords{示向度交会\quad 定位误差\quad 信息驱动搜索\quad 路径规划\quad 定向干扰源}

 \end{abstract}

 % ---------------------- 正文（不要目录，不超过 30 页） ----------------------
 % =====================================================================
%  一、问题重述
% =====================================================================
\section{问题重述}

\subsection{问题背景}
无线电测向与干扰源定位是频谱管理、电磁环境监测中的重要任务。借助可移动测向机，
通过在不同检测点测得干扰源的示向度（方位角），即可对干扰源位置进行交会估计；
在此基础上进一步规划检测路线、逐一清除干扰源。\textbf{【待填充：结合赛题原文，
补充背景与任务场景，注意不要照抄原题，用自己语言概述。】}

\subsection{问题的提出}
根据赛题，本文需要依次解决以下四个子问题：

\begin{enumerate}
    \item \textbf{问题一}：……在已知两个检测点及其示向度（含测量误差）的情况下，
          求交会定位区域的直径，并判断以该直径为直径的圆能否覆盖整个定位区域。
          \textbf{【待填充：按原题准确复述。】}

    \item \textbf{问题二}：……在给定第一检测点测得示向度后，研究第二检测点的
          最优选择策略，并给出相应的候选区域。
          \textbf{【待填充：按原题准确复述。】}

    \item \textbf{问题三}：……在半径为 $1800\,\mathrm{m}$ 的圆形区域内存有
          $10$--$16$ 个位置未知的全向干扰源，要求在规定现实时间（$20$ 分钟）内
          将全部干扰源定位并清除。
          \textbf{【待填充：按原题准确复述，含测向、换频道、移动、清除等时间约束。】}

    \item \textbf{问题四}：……在问题三基础上叠加若干方向未知的定向干扰源，
          要求仍能在限定时间内完成定位与清除。
          \textbf{【待填充：按原题准确复述。】}
\end{enumerate}
     % 一、问题重述
 % =====================================================================
%  二、问题分析
% =====================================================================
\section{问题分析}

\subsection{总体思路}
四个子问题层层递进：问题一建立最基本的示向度交会几何定位模型；问题二将其
扩展为检测点的最优实验设计；问题三在前述基础上，面对多个未知源与严格的时间
预算，形成信息驱动的大范围搜索与清除策略；问题四进一步处理定向干扰源带来的
方向不确定性与诊断困难。\textbf{【待填充：补充总体技术路线，可用一张流程图
（图~1）概括。】}

\subsection{问题一分析}
问题一的本质是\textbf{计算几何问题}：每条示向度带有 $\pm \delta$ 的误差，形成
一条扇形射线束，两条射线束的交叠区域即定位区域；区域直径可由凸包直径
（旋转卡壳算法）求得，而“直径圆能否覆盖该区域”则归结为最小覆盖圆问题。
\textbf{【待填充：细化分析，说明“直径圆可能无法覆盖细长区域”的直觉。】}

\subsection{问题二分析}
问题二具有\textbf{最优实验设计}的思想：示向度误差 $\delta$ 在距离 $d$ 处产生的
横向误差约为 $d\tan\delta$，故“离目标近”与“两条示向度交会角接近 $90^\circ$”
之间存在权衡，可用 Fisher 信息矩阵 / GDOP / 误差椭圆加以形式化。
\textbf{【待填充：给出第二检测点选择的优化准则与候选区域的刻画方法。】}

\subsection{问题三分析}
问题三是本题的主战场，受\textbf{严格现实时间预算}（$20$ 分钟）约束。仅靠
“边扫边定”的低效策略难以完成，需要利用“频道数量有限且互不相同”这一结构，
先粗扫建立“频道—射线束”台账，再对交会区域收敛的频道精确定位，本质上是一个
\textbf{信息驱动的旅行商问题（TSP with active sensing）}。
\textbf{【待填充：展开时间预算核算与整体策略框架。】}

\subsection{问题四分析}
问题四叠加了定向干扰源（方向未知），使“测不到信号”出现了多种截然不同的
原因（无该频道、超出有效接收半径、处于定向覆盖角之外），诊断逻辑复杂度显著
上升。\textbf{【待填充：分析定向源对定位与清除流程的影响与应对思路。】}
    % 二、问题分析
 % =====================================================================
%  三、模型假设
% =====================================================================
\section{模型假设}
为便于建模与求解，在合理范围内作出如下假设：

\begin{enumerate}
    \item 测向机的示向度测量误差可视为零均值、有界的随机误差，其界为赛题给定的
          $\pm\delta$；\textbf{【待填充：按实际情况修订。】}
    \item 问题三中的干扰源均为\textbf{全向辐射源}，各向辐射特性相同；
    \item 各干扰源位置在检测与清除过程中\textbf{固定不变}，不随时间移动；
    \item 测向机以恒定速度移动（$5\,\mathrm{m/s}$），忽略转向、加减速耗时；
    \item 不同干扰源的工作频道\textbf{互不相同}，同一频道只需定位一次；
    \item 检测区域内无障碍遮挡，测向机与干扰源之间可视；
    \item \textbf{【待填充：可根据模型需要增删假设，并说明每条假设的合理性。】}
\end{enumerate}
 % 三、模型假设
 % =====================================================================
%  四、符号说明
% =====================================================================
\section{符号说明}
本文使用的主要符号及其含义见表~\ref{tab:symbol}。

\begin{table}[!htbp]
    \caption{主要符号说明}\label{tab:symbol}
    \centering
    \begin{tabular}{ll}
        \toprule
        符号 & 含义 \\
        \midrule
        $\theta_i$      & 第 $i$ 次测向得到的示向度（方位角） \\
        $\delta$        & 示向度测量误差（$^\circ$） \\
        $d$             & 检测点与干扰源之间的距离（m） \\
        $R$             & 干扰源有效接收半径（m） \\
        $P_i$           & 第 $i$ 个检测点坐标 \\
        $S_j$           & 第 $j$ 个干扰源坐标（待估计量） \\
        $v$             & 测向机移动速度（$\mathrm{m/s}$） \\
        $T$             & 现实时间预算（s） \\
        $\varepsilon$   & 定位误差 / 置信椭圆半轴参数 \\
        \bottomrule
    \end{tabular}
\end{table}

\textbf{【待填充：符号说明须与后文模型中的符号一一对应，可按需增删行；未在此
列出的符号在首次出现处说明。】}
     % 四、符号说明
 % =====================================================================
%  五、模型建立与求解（核心章节）
%  每个子问题统一采用「建模思路 → 模型建立 → 算法/求解 → 结果分析」结构。
% =====================================================================
\section{模型建立与求解}

% ---------------------------------------------------------------------
\subsection{问题一：示向度交会定位模型}
\subsubsection{建模思路}
\textbf{【待填充：说明由两条含误差的示向度交会出定位区域的基本思想。】}

\subsubsection{模型建立}
设检测点 $P_1(x_1,y_1)$、$P_2(x_2,y_2)$ 处测得示向度分别为 $\theta_1$、$\theta_2$，
考虑测量误差 $\delta$，则真实方位落在 $[\theta_i-\delta,\theta_i+\delta]$ 内，两条
扇形射线束的交叠区域即为干扰源的可能位置区域。
\textbf{【待填充：写出射线/扇形边界的解析表达式，定义交会定位区域。】}

区域\textbf{直径}定义为区域内任意两点间距离的上确界，等价于凸包的直径，可用
旋转卡壳（Rotating Calipers）算法在 $O(n\log n)$ 内求得；“直径圆”能否覆盖区域
则归结为\textbf{最小覆盖圆}（smallest enclosing circle）问题。
\textbf{【待填充：给出直径与最小覆盖圆的数学定义及判据。】}

\subsubsection{求解与结果}
\textbf{【待填充：给出问题一的求解流程、示例交会区域示意图（图~2）与结论。】}

% ---------------------------------------------------------------------
\subsection{问题二：第二检测点最优选择模型}
\subsubsection{建模思路}
\textbf{【待填充：说明“距离近”与“交会角接近 $90^\circ$”之间的权衡。】}

\subsubsection{模型建立}
在距离 $d$ 处，示向度误差 $\delta$ 产生的横向定位误差约为 $d\tan\delta$。两条
示向度交会时，可用\textbf{Fisher 信息矩阵}刻画定位不确定度，其行列式
$\det(\mathbf{F})$ 与交会角及两检测点到目标的距离有关；等价地可用 GDOP 或
误差椭圆面积作为准则。候选区域可由等精度线 / 扇环刻画。
\textbf{【待填充：写出 Fisher 信息矩阵 / 误差椭圆的具体形式与最优化目标。】}

\subsubsection{求解与结果}
\textbf{【待填充：给出第二检测点选择策略、候选区域图示（图~3）与结果。】}

% ---------------------------------------------------------------------
\subsection{问题三：多源快速定位与清除策略}
\subsubsection{建模思路}
\textbf{【待填充：说明“先粗扫建频道台账，再精定位收敛频道”的信息驱动策略，以及
时间预算的核算（移动、测向、换频道、清除各自耗时）。】}

\subsubsection{模型建立}
将测向机视为在 $1800\,\mathrm{m}$ 半径圆域内移动的智能体，待清除的 $10$--$16$
个干扰源位置未知。该问题可建模为\textbf{信息驱动的旅行商问题（TSP with active
sensing）}：以栅格或正六边形铺排搜索路线，用射线束交点与置信椭圆修剪候选区域，
依据“频道—射线束”台账决定下一次测向位置与频道。
\textbf{【待填充：给出目标函数（如总清除时间最小 / 清除率最大）、状态转移与
决策规则。】}

\subsubsection{求解与结果}
\textbf{【待填充：给出算法流程（可配伪代码）、模拟器联调结果、清除比例与平均
定位清除时间等实测指标，附表格（表~2）与轨迹图（图~4）。】}

% ---------------------------------------------------------------------
\subsection{问题四：定向干扰源定位模型}
\subsubsection{建模思路}
\textbf{【待填充：说明定向源方向未知带来的“测不到信号”多重成因与诊断逻辑。】}

\subsubsection{模型建立}
定向干扰源仅在有限方位覆盖角内可被检测，故“未测得信号”需区分三种原因：该频道
无干扰源、超出有效接收半径、处于定向覆盖角之外。建模需在问题三的搜索框架上引入
\textbf{定向覆盖角与源朝向}两个未知量，并设计相应的绕测 / 诊断策略。
\textbf{【待填充：给出定向源的可检测条件、参数估计模型与诊断规则。】}

\subsubsection{求解与结果}
\textbf{【待填充：给出问题四的求解流程与结果，并与问题三对比。】}
       % 五、模型建立与求解
 % =====================================================================
%  六、模型评价与改进
% =====================================================================
\section{模型评价与改进}

\subsection{灵敏度分析}
\textbf{【待填充：对示向度误差 $\delta$、移动速度 $v$、有效接收半径 $R$ 等关键参数
做灵敏度分析，讨论其对定位精度与清除时间的影响，可配图（图~5）。】}

\subsection{模型优点}
\begin{itemize}
    \item \textbf{【待填充：例如——将定位问题形式化为计算几何与最优实验设计，理论清晰。】}
    \item \textbf{【待填充：……】}
\end{itemize}

\subsection{模型缺点}
\begin{itemize}
    \item \textbf{【待填充：例如——问题三/四高度依赖模拟器协议与在线联调，存在不确定性。】}
    \item \textbf{【待填充：……】}
\end{itemize}

\subsection{改进方向}
\textbf{【待填充：例如——引入自适应步长的贝叶斯搜索、在线学习策略、更精细的
时间—精度权衡等。】}
  % 六、模型评价与改进
 % =====================================================================
%  七、参考文献
%  注：cumcmthesis 已将 thebibliography 环境重定义为自带「参考文献」标题，
%      故此处不要再写 \section{参考文献}，避免重复标题。
% =====================================================================
\begin{thebibliography}{9}
    \bibitem{bib:1} 作者. 无线电测向与定位技术[M]. 出版社, 年份.
    \bibitem{bib:2} 作者. 论文题目[J]. 期刊名, 年份, 卷(期): 页码.
    \bibitem{bib:3} 全国大学生数学建模竞赛组委会. 全国大学生数学建模竞赛论文格式规范（2026年修订稿）.
    % 【待填充：按科技论文规范列出参考文献，并在正文引用处用 \cite{...} 标注。】
\end{thebibliography}
  % 七、参考文献

 % ---------------------- 附录（页数不限） ----------------------
 % =====================================================================
%  附录（页数不限）
%  依据 2026 规范：附录必须包含支撑材料的文件列表、建模所用全部完整可运行的
%  源程序代码；如确无程序须注明“本论文没有用到程序”。
% =====================================================================
\begin{appendices}

\section{支撑材料文件列表}
本论文支撑材料所包含的文件见表~\ref{tab:files}（对应压缩包内的相对路径）。

\begin{table}[!htbp]
    \caption{支撑材料文件列表}\label{tab:files}
    \centering
    \begin{tabular}{cll}
        \toprule
        序号 & 文件路径 & 说明 \\
        \midrule
        1 & \texttt{code/main.py}      & 主程序（测向与搜索策略） \\
        2 & \texttt{code/}…            & 【待填充：其余源程序】 \\
        3 & \texttt{data/}…            & 【待填充：自主查阅使用的数据资料】 \\
        4 & \texttt{figures/}…         & 【待填充：较大篇幅中间结果图表】 \\
        \bottomrule
    \end{tabular}
\end{table}

\textbf{【待填充：最终按实际支撑材料逐项填写；如确无支撑材料，则注明“本论文没有
支撑材料”。】}

\section{主要源程序}
\textbf{【待填充：放入建模所用到的全部完整、可运行的源程序代码（含 EXCEL、SPSS
等交互命令）。下方为占位示例。】}

\begin{lstlisting}[language=python, caption={测向与搜索策略主程序（占位示例）}]
# -*- coding: utf-8 -*-
# 占位示例：B 题测向机在线搜索与清除策略的主流程骨架
def main():
    # TODO: 建立到模拟器的连接，发送 /enter 进入任务
    # TODO: 粗扫建立「频道 -> 射线束」台账
    # TODO: 对交会区域收敛的频道进行精确定位
    # TODO: 逐个 /clear 清除干扰源，并在时间预算内返回
    pass

if __name__ == "__main__":
    main()
\end{lstlisting}

\end{appendices}


\end{document}
